\documentclass[12pt]{article}
\usepackage[utf8]{inputenc}
\usepackage[T1]{fontenc}
\usepackage[spanish]{babel}
\usepackage{graphicx}
\usepackage{tikz}
\usepackage{float}
\usepackage{enumitem}
\usepackage[table,xcdraw]{xcolor}
\usepackage{array}
\usepackage{caption}
\usepackage[square,numbers]{natbib}
\usepackage{amsmath,amssymb}
\usepackage[absolute,overlay]{textpos}
\usepackage{setspace}
\usepackage{booktabs}
\usepackage{tabularx}
\usepackage{listings}
\usepackage{subcaption}
\usepackage[square,numbers]{natbib}
\decimalpoint

\usepackage[hyphens]{url}
\usepackage[hidelinks]{hyperref}
\usepackage{breakurl}
\urlstyle{same}   

\usepackage{ebgaramond}   % Fuente del contenido
\usepackage{helvet}       % Para títulos de secciones
\renewcommand{\familydefault}{\rmdefault}
\usepackage{sectsty}
\allsectionsfont{\sffamily\bfseries} 

\newcommand{\TitleStyleOne}[1]{\bfseries\sffamily #1}

\definecolor{codegreen}{rgb}{0,0.6,0}
\definecolor{codegray}{rgb}{0.5,0.5,0.5}
\definecolor{codepurple}{rgb}{0.58,0,0.82}
\definecolor{backcolour}{rgb}{0.96,0.96,0.96}

\lstdefinestyle{pystyle}{
    backgroundcolor=\color{backcolour},   
    commentstyle=\color{codegreen},
    keywordstyle=\color{magenta},
    numberstyle=\tiny\color{codegray},
    stringstyle=\color{codepurple},
    basicstyle=\ttfamily\scriptsize,
    breakatwhitespace=false,         
    breaklines=true,                 
    captionpos=b,                    
    keepspaces=true,                 
    numbers=left,                    
    numbersep=5pt,                  
    showspaces=false,                
    showstringspaces=false,
    showtabs=false,                  
    tabsize=2
}
\lstset{style=pystyle}

\begin{document}

\begin{textblock*}{9cm}(\dimexpr0.5\paperwidth-4.5cm\relax,2cm)
\color{black}
\centering
\footnotesize\textbf{Universidad Autónoma de Nuevo León}\\
\footnotesize Maestría en Ciencia de Datos\\
\footnotesize Procesamiento y Clasificación de Datos
\end{textblock*}

\vspace*{0cm}

\begin{center}
    {\large \TitleStyleOne{Análisis Multimodal y Clasificación de Datos No Estructurados}}\\[0.5em]
    
    % Línea fina debajo del título
    \rule{0.65\textwidth}{0.35pt}\\[0.7em]
    
    % Datos
    {\large{Jessica Lizeth Hernández Bracho -- 1842553}}\\[0.3em]
    {\today}
\end{center}

\vspace{0.8cm} 

\section{Introducción}

Aproximadamente el 80\% de los datos generados a nivel global corresponden a formatos no estructurados como texto libre, imágenes digitales y señales acústicas. A diferencia de las estructuras tabulares convencionales, la información no estructurada requiere técnicas avanzadas de normalización, extracción de características y modelado no lineal para su correcta interpretación computacional.

El presente trabajo integra la metodología y los hallazgos analíticos desarrollados a lo largo del curso, abordando el procesamiento de datos no estructurados en tres dimensiones fundamentales: Análisis Léxico y Semántico de Texto, Visión por Computadora (segmentación industrial y \textit{Transfer Learning} biomédico) y Procesamiento de Señales de Audio y Voz (dominio frecuencial, descomposición Wavelet y alineación temporal DTW).

\section{Descripción de los Datos}
El análisis se sustentó en un conjunto de datos organizado en tres ejes temáticos:

\begin{itemize}
    \item \textbf{Textual de Series Televisivas:} Guiones en texto plano (\texttt{.txt}) de los episodios piloto ($1\times01$) de las comedias \textit{Malcolm in the Middle} (Linwood Boomer) y \textit{Modern Family} (Christopher Lloyd y Steven Levitan).
    \item \textbf{Textual Académico:} Subconjunto balanceado del repositorio oficial \textit{20 Newsgroups} (\texttt{sklearn.datasets.fetch\_20newsgroups}), aislando las categorías competitivas, para un total de 1,892 observaciones.
    \item \textbf{Visión Industrial y Biomédica:} Fotografía industrial de balines de acero y dataset biomédico \textit{Malaria} (500 micrografías de eritrocitos: 400 de entrenamiento y 100 de validación).
    \item \textbf{Audio y Voz:} Intros musicales en formato MP3 (\textit{Game of Thrones} vs \textit{The Big Bang Theory}) y dataset \textit{Spoken MNIST} compuesto por señales de voz aisladas muestreadas a $8\text{ kHz}$ (dígitos $0$ y $1$).
\end{itemize}

\section{Antecedentes}
El tratamiento de datos no estructurados ha evolucionado desde representaciones heurísticas hacia modelos vectoriales densos y redes profundas. En minería de texto, las métricas de riqueza léxica (TTR), modelos supervisados de polaridad (\textit{AFINN-165}) y la vectorización TF-IDF con lematización basada en \textit{WordNet} fundamentan la clasificación automática. En visión por computadora, las transformaciones geométricas y los detectores de bordes (Canny, Hough) permiten la metrología industrial automatizada, mientras que el aprendizaje por transferencia (\textit{Transfer Learning}) con arquitecturas como \textit{VGG16} reutiliza representaciones jerárquicas aprendidas en ImageNet. Finalmente, en señales acústicas, la caracterización espectral (STFT, MFCC, Centroides) y la descomposición multirresolución mediante transformada Wavelet acoplada a \textit{Dynamic Time Warping} (DTW) permiten clasificar secuencias de voz con variaciones temporales no lineales.

\section{Metodología}

\subsection{Módulo de Texto: Procesamiento, Estilo y Clasificación}

El procesamiento y análisis del lenguaje natural se estructuró en dos vertientes principales: la caracterización estilístico-emocional de guiones de comedia (\textit{Malcolm in the Middle} y \textit{Modern Family}) y la clasificación temática supervisada sobre la identificación dentro de \textit{20 Newsgroups}.

\subsubsection{Preprocesamiento y Limpieza Adaptativa de Guiones}
Los guiones de los episodios piloto en texto plano (\texttt{.txt}) sufrieron una transformación en varias etapas para eliminar ruido textual sin perder la esencia sintáctica del diálogo:
\begin{enumerate}
    \item \textbf{Filtrado de Acotaciones Escénicas:} Mediante expresiones regulares (\texttt{re.sub}), se removieron las anotaciones entre paréntesis y corchetes que corresponden a instrucciones del director y no al diálogo hablado.
    \item \textbf{Normalización Estándar:} Conversión total a minúsculas, eliminación de signos de puntuación y caracteres especiales no alfabéticos.
    \item \textbf{Remoción Extendida de Stopwords:} Se utilizó la lista base en inglés del paquete \texttt{nltk.corpus} complementada con jerga conversacional propia de guiones de televisión (\texttt{gonna}, \texttt{wanna}, \texttt{gotta}, \texttt{yeah}, \texttt{oh}).
    \item \textbf{Reducción Morfológica (Stemming):} Se aplicó el algoritmo \texttt{PorterStemmer} para reducir los tokens a su raíz morfológica común.
\end{enumerate}

\begin{lstlisting}[language=Python, caption={Preprocesamiento y limpieza adaptativa para guiones de series.}]
import re
import nltk
from nltk.corpus import stopwords
from nltk.stem import PorterStemmer

def pipeline_limpieza_guion(texto_raw):
    # 1. Eliminar acotaciones escenicas [texto] o (texto)
    texto_limpio = re.sub(r'\[.*?\]|\(.*?\)', '', texto_raw)
    
    # 2. Normalizacion a minusculas y caracteres alfabeticos
    texto_limpio = re.sub(r'[^a-zA-Z\s]', '', texto_limpio).lower()
    
    # 3. Tokenizacion y filtrado de stopwords extendidas
    tokens = texto_limpio.split()
    stop_words_ext = set(stopwords.words('english')).union({'gonna', 'wanna', 'gotta', 'yeah', 'oh'})
    tokens_filtrados = [w for w in tokens if w not in stop_words_ext]
    
    # 4. Stemming con Porter
    stemmer = PorterStemmer()
    tokens_stemmed = [stemmer.stem(w) for w in tokens_filtrados]
    
    return tokens_filtrados, tokens_stemmed
\end{lstlisting}

\subsubsection{Análisis Estilístico, Dinámica Narrativa y Sentimiento}
Para capturar la evolución de la narrativa, cada guión se dividió cronológicamente en 4 bloques homogéneos (correspondientes a los actos dramáticos). Se evaluó la \textbf{Riqueza Léxica} mediante la métrica \textit{Type-Token Ratio} ($\text{TTR}$):

$$\text{TTR} = \frac{|V|}{N}$$

donde $|V|$ representa el número de palabras únicas (vocabulario) y $N$ el total de tokens en el bloque.

La \textbf{similitud léxica} entre ambas producciones se cuantificó mediante la Similitud Coseno sobre la matriz cruzada de frecuencias de términos:

$$\text{Similitud}(A, B) = \cos(\theta) = \frac{\mathbf{A} \cdot \mathbf{B}}{\|\mathbf{A}\| \|\mathbf{B}\|} = \frac{\sum_{i=1}^{n} A_i B_i}{\sqrt{\sum_{i=1}^{n} A_i^2} \sqrt{\sum_{i=1}^{n} B_i^2}}$$

Para el \textbf{análisis de sentimiento} se mapearon los tokens filtrados contra el léxico supervisado \textit{AFINN-165}, el cual asigna una puntuación entera entre $-5$ (emoción extremadamente negativa) y $+5$ (extremadamente positiva). La puntuación de sentimiento promedio por acto se definió como:

$$\bar{S} = \frac{1}{|T_{\text{AFINN}}|} \sum_{w \in T_{\text{AFINN}}} \text{Score}(w)$$

\subsubsection{Lematización y Clasificación Temática en 20 Newsgroups}
Para el análisis de \textit{20 Newsgroups} (\texttt{rec.autos} vs. \texttt{sci.space}), se reemplazó el \textit{stemming} por una \textbf{lematización morfológica} basada en la base de datos léxica \textit{WordNet} (\texttt{WordNetLemmatizer}), preservando la validez gramatical de los términos.

\begin{table}[H]
\resizebox{\textwidth}{!}{%
\begin{tabular}{lcccc}
\hline
\rowcolor[HTML]{CBCEFB} 
\textbf{Categoría} & \textbf{Docs} & \textbf{Prom. Palabras (Std)} & \textbf{Prom. Caracteres (Std)} & \textbf{Densidad Léxica} \\ \hline
rec.autos          & 937           & 119.53 (221.18)               & 692.21 (1342.46)                & 5.61                     \\
sci.space          & 955           & 196.27 (502.94)               & 1253.26 (3367.81)               & 5.64                     \\ \hline
\end{tabular}%
}
\caption{Distribución de Clases (Autos vs Espacio)}
\end{table}

Con el fin de determinar la mejor representación vectorial se diseñó un experimento comparativo evaluando 4 esquemas de extracción de características:
\begin{itemize}
    \item \textbf{Conteo Simple (Bag of Words):} Matriz de frecuencia absoluta de términos ($tf_{t,d}$).
    \item \textbf{TF-IDF a Nivel Palabra:} Ponderación por Frecuencia de Término e Inversa de Frecuencia de Documento, como se muestra en la siguiente formula
    $$\text{TF-IDF}(t, d, D) = tf_{t,d} \times \log \left( \frac{|D|}{|\{d \in D : t \in d\}|} \right)$$
    \item \textbf{TF-IDF a Nivel N-Gramas:} Extracción de secuencias contiguas de palabras en un rango de $(2, 3)$-gramas para capturar contexto local.
    \item \textbf{TF-IDF a Nivel Carácter:} Factorización en subpalabras y n-gramas de caracteres ($n \in [3, 5]$) para otorgar robustez ante errores ortográficos y variaciones morfológicas.
\end{itemize}

Cada representación vectorial se evaluó mediante un clasificador de \textbf{regresión logística}, optimizando la función de pérdida con regularización $L_2$ e iterando sobre tres valores del hiperparámetro de fuerza de regularización $C \in \{0.1, 1.0, 10.0\}$.

\subsection{Módulo de Visión Industrial: Transformada de Hough}
Para la inspección automatizada de balines de acero, la imagen de $900 \times 900$ píxeles se convirtió a escala de grises y se le aplicó un Suavizado Gaussiano con un kernel de $(13, 13)$ para atenuar los reflejos especulares. La extracción de fronteras y la localización de centros geométricos se realizó mediante el algoritmo \texttt{cv2.HoughCircles}:

\begin{lstlisting}[language=Python, caption={Procesamiento para detección de balines.}]
import cv2
import numpy as np

gray = cv2.cvtColor(img, cv2.COLOR_BGR2GRAY)
blurred = cv2.GaussianBlur(gray, (13, 13), 0)
edges = cv2.Canny(blurred, 30, 150)

circles = cv2.HoughCircles(
    blurred,
    cv2.HOUGH_GRADIENT,
    dp=1.2,
    minDist=35,
    param1=50,
    param2=25,
    minRadius=10,
    maxRadius=25
)
\end{lstlisting}

\subsection{Módulo de Visión Biomédica: Transfer Learning con VGG16}
Para la detección de eritrocitos infectados con \textit{Malaria}, las micrografías se reescalaron a $(150, 150, 3)$ y se normalizaron en el rango $[0, 1]$. Se utilizó la arquitectura \textbf{VGG16} preentrenada en ImageNet congelando todos sus bloques convolucionales (\texttt{trainable=False}) y acoplando una cabeza de clasificación densa:

\begin{lstlisting}[language=Python, caption={Construcción y compilación de la red VGG16.}]
from tensorflow.keras.applications import VGG16
from tensorflow.keras import layers, models

base_model = VGG16(weights='imagenet', include_top=False, input_shape=(150, 150, 3))
base_model.trainable = False

model = models.Sequential([
    base_model,
    layers.Flatten(),
    layers.Dense(256, activation='relu'),
    layers.Dense(1, activation='sigmoid')
])

model.compile(optimizer='adam', loss='binary_crossentropy', metrics=['accuracy'])
\end{lstlisting}

\subsection{Módulo de Audio y Voz: Análisis Frecuencial, Wavelet y DTW}
Se extrajo los coeficientes MFCC (13 componentes) y la STFT para cuantificar el \textbf{Centroide Espectral Medio} ($\mu_{\text{cent}}$) de intros musicales. 

Así como, para la clasificación de dígitos hablados (\textit{Spoken MNIST}), las señales se normalizaron primero en amplitud dividiendo entre la magnitud máxima para anular diferencias de volumen ($\|a\|_{\infty} = 1$). Posteriormente, se aplicó una descomposición Wavelet discreta Daubechies 1 (\texttt{db1}) a 4 niveles de resolución, extrayendo el vector de coeficientes de aproximación $cA_4$ como la firma acústica de baja frecuencia. La distancia entre la señal incógnita y las referencias conocidas ($0$ y $1$) se calculó mediante \textbf{Dynamic Time Warping (DTW)}:

\begin{lstlisting}[language=Python, caption={Normalización, descomposición Wavelet y clasificación DTW.}]
import pywt
from dtw import dtw

# Normalización de amplitud
a_A_norm = audio_A / np.max(np.abs(audio_A))
a_B_norm = audio_B / np.max(np.abs(audio_B))
a_test_norm = audio_test / np.max(np.abs(audio_test))

# Descomposición Wavelet (Nivel 4, db1)
coeffs_A = pywt.wavedec(a_A_norm, 'db1', level=4)[0]
coeffs_B = pywt.wavedec(a_B_norm, 'db1', level=4)[0]
coeffs_test = pywt.wavedec(a_test_norm, 'db1', level=4)[0]

# Clasificación por mínima distancia acumulada DTW
res_A = dtw(coeffs_test, coeffs_A, keep_internals=True)
res_B = dtw(coeffs_test, coeffs_B, keep_internals=True)
\end{lstlisting}

\section{Resultados y Discusión}

\subsection{Análisis Léxico y Emocional de Guiones}
Las palabras más frecuentes en \textit{Malcolm} correspondieron a \texttt{malcolm} (33), \texttt{okay} (25), \texttt{know} (24) y \texttt{go} (23); mientras que en \textit{Modern Family} destacaron \texttt{know} (33), \texttt{okay} (30), \texttt{that} (24) y \texttt{dad} (21).

\begin{table}[H]
\resizebox{\textwidth}{!}{%
\begin{tabular}{lcc}
\hline
\rowcolor[HTML]{CBCEFB} 
\textbf{Métrica}                & \textbf{Malcolm in the Middle} & \textbf{Modern Family}  \\ \hline
Signos de Exclamación (!)       & 55                             & 62                      \\
Signos de Interrogación (?)     & 89                             & 104                     \\
Riqueza Léxica                  & \textbf{0.423 (42.3\%)}        & \textbf{0.400 (40.0\%)} \\
Carga Positiva                  & \textbf{67.57\%}               & 59.84\%                 \\
Carga Negativa                  & 32.43\%                        & 40.16\%                 \\
Puntuación Promedio Sentimiento & \textbf{0.605}                 & 0.389                   \\
Top 1 Bigrama                   & (wait, wait) {[}5{]}           & (hi, hi) {[}4{]}        \\
Top 2 Bigrama                   & (hold, hand) {[}4{]}           & (thank, thank) {[}4{]}  \\
Top 3 Bigrama                   & (special, class) {[}3{]}       & (let, see) {[}4{]}      \\ \hline
\end{tabular}%
}
\caption{Métricas estilísticas, léxicas y emocionales}
\end{table}

El análisis por bloques temporales mostró que \textit{Malcolm} se mantiene estable entre 210 y 230 palabras únicas por acto como se puede ver en la Figura \ref{fig:series}, mientras que \textit{Modern Family} alcanza picos superiores a 260 palabras únicas en los actos intermedios. La distancia de Coseno fue de \textbf{0.3586}, registrando una similitud léxica del \textbf{64.14\%}.

\begin{figure}[H]
    \centering
    \includegraphics[width=0.72\textwidth]{Imagenes/graf1(serie).png}
    \caption{Diversidad léxica por bloques.}
    \label{fig:series}
\end{figure}

\subsection{Clasificación de Texto en 20 Newsgroups}
El modelo \textbf{TF-IDF a Nivel Palabra} con $C=1.0$ alcanzó el desempeño superior con una exactitud global de \textbf{94.72\%}.

\begin{table}[H]
\resizebox{\textwidth}{!}{%
\begin{tabular}{lcccl}
\hline
\rowcolor[HTML]{CBCEFB} 
\textbf{Estrategia Vectorial} & \textbf{C=0.1} & \textbf{C=1.0}   & \textbf{C=10.0} & \textbf{Diagnóstico Técnico}         \\ \hline
Conteo Simple         & 94.20\% & 94.20\% & 91.03\% & Sobreajuste a alta C por palabras ruidosas.   \\
\textbf{TF-IDF Palabras}      & 91.20\%        & \textbf{94.72\%} & 94.46\%         & \textbf{Punto de equilibrio óptimo.} \\
TF-IDF N-Gramas (2,3) & 53.30\% & 87.34\% & 88.39\% & Subajuste severo en C=0.1 por dispersión.     \\
TF-IDF Caracteres     & 67.28\% & 69.80\% & 70.71\% & Baja capacidad discriminativa en texto largo. \\ \hline
\end{tabular}%
}
\caption{Matriz de exactitud del diseño de experimentos en 20 Newsgroups.}
\end{table}

\subsection{Metrología Industrial con OpenCV}
El algoritmo Hough detectó de forma automatizada un total de \textbf{134 balines} en la muestra industrial. El rango de radios estimados osciló entre $12\text{ px}$ y $24\text{ px}$, logrando aislar correctamente los centroides geométricos $(\bar{X}, \bar{Y})$ incluso en presencia de variaciones por iluminación reflexiva.

\begin{figure}[H]
    \centering
    \includegraphics[width=0.8\textwidth]{Imagenes/ball1.png}
    \caption{Detección de Bordes (Canny) y Detección de Círculos (Hough).}
    \label{fig:opencv1}
\end{figure}

\subsection{Clasificación Biomédica con VGG16}
El conjunto de datos seleccionado para este experimento es el dataset oficial Malaria, disponible de manera nativa a través del repositorio de la API de \textit{tensorflow datasets (TFDS)}.

El conjunto de datos global está compuesto por $27{,}558$ imágenes individuales de células sanguíneas, distribuidas equitativamente en dos categorías biológicas:

\begin{itemize}
    \item \textbf{Clase 0 --- \textit{Infectada}:} Células que albergan al parásito \textit{Plasmodium falciparum} en alguna de sus fases evolutivas. Morfológicamente se caracterizan por la presencia de pequeños anillos densos o inclusiones punctiformes de tonalidad morada/azul brillante dentro del citoplasma, derivadas de la tinción de Giemsa.
    \item \textbf{Clase 1 --- \textit{Sana / Normal}}: Eritrocitos limpios que exhiben una textura y coloración rosada/rojiza homogénea, carentes de estructuras o cuerpos extraños intracelulares.
\end{itemize}

A partir del universo total, se extrajo un subconjunto estricto y equilibrado para garantizar la evaluación computacional del modelo:

\begin{itemize}
    \item \textbf{Subconjunto de Entrenamiento:} 400 micrografías digitales.
    \item \textbf{Subconjunto de Validación:} 100 micrografías digitales.
\end{itemize}

La arquitectura convolucional acumuló $14,714,688$ parámetros no entrenables (base VGG16) y $2,097,665$ parámetros entrenables en el clasificador denso. Tras 10 épocas, la red logró una exactitud en entrenamiento de \textbf{94.25\%} con una pérdida de $0.1633$. En el conjunto de prueba independiente (100 imágenes), la exactitud de validación se estabilizó en \textbf{77.00\%} con un \textit{val\_loss} de $0.4343$.

\begin{figure}[H]
    \centering
    \includegraphics[width=0.8\textwidth]{Imagenes/results(t5).png}
    \caption{Gráficas de Resultados.}
    \label{fig:biomed}
\end{figure}

\subsection{Análisis Acústico de Intros Musicales}
Se seleccionaron como casos de estudio dos series con estructuras de producción opuestas: el tema principal de Game of Thrones (representativo del género dramático/orquestal) y un intro típico de The Big Bang Theory (representativo del género de comedia/pop).

\begin{figure}[H]
    \centering
    \includegraphics[width=0.8\textwidth]{Imagenes/onda1(series).png}
    \caption{Comparativo de formas de onda.}
    \label{fig:Adseries}
\end{figure}

El análisis frecuencial mediante STFT reveló diferencias estructurales significativas entre géneros de apertura:
\begin{itemize}
    \item \textbf{Game of Thrones:} Centroide Espectral Medio de \textbf{$1,325.89$} Hz (timbre grave/oscuro).
    \item \textbf{The Big Bang Theory:} Centroide Espectral Medio de \textbf{$2,622.35$} Hz (timbre brillante/agudo).
    \item \textbf{Brecha Espectral Absoluta:} \textbf{$1,296.46$} Hz.
\end{itemize}

\begin{figure}[H]
    \caption{Frecuencias Representativas - Intro de Series.}
    \centering
    \begin{subfigure}[b]{0.7\textwidth}
        \centering
        \includegraphics[width=\textwidth]{Imagenes/frGOT(series).png}
        \caption{Game of Thrones}
        \label{subfig:got}
    \end{subfigure}
    \hfill
    \begin{subfigure}[b]{0.7\textwidth}
        \centering
        \includegraphics[width=\textwidth]{Imagenes/frTBBT(series).png}
        \caption{The Big Bang Theory}
        \label{subfig:tbbt}
    \end{subfigure}
    
    \label{fig:frec.rep}
\end{figure}

\subsection{Reconocimiento de Voz y Alineación DTW}
A diferencia de la Transformada de Fourier, la Transformada Wavelet Discreta empleando la familia Daubechies (\texttt{db1} / Haar) permite descomponer la señal de audio en múltiples niveles de resolución temporal y frecuencial:

\begin{itemize}
    \item \textbf{Coeficientes de Aproximación ($cA$):} Capturan las frecuencias bajas y la envolvente fonética principal de la voz (es decir, el contorno espectral al articular palabras como \textit{zero} / \textit{one}).
    \item \textbf{Coeficientes de Detalle ($cD$):} Capturan las frecuencias altas y variaciones de rápida transición (asociadas a ruido de fondo o consonantes sibilantes).
\end{itemize}

Para esta prueba, se seleccionaron los coeficientes de aproximación del cuarto nivel ($cA_4$), logrando una reducción significativa en la dimensión temporal del audio sin comprometer la estructura morfofonética subyacente.

\begin{figure}[H]
    \centering
    \includegraphics[width=0.8\textwidth]{Imagenes/onda2(dig).png}
    \caption{Visualización de las Formas de Onda.}
    \label{fig:AdVoz}
\end{figure}

Para la muestra de prueba correspondiente al dígito real "$0$", arrojó los siguientes resultados:

\begin{itemize}
    \item \textbf{Distancia acumulada hacia Muestra A (Dígito '$0$'):} $70.7006$
    \item \textbf{Distancia acumulada hacia Muestra B (Dígito '$1$'):} $75.4533$
\end{itemize}

El criterio del vecino más cercano ($k\text{-NN}, k=1$) seleccionó correctamente la \textbf{Muestra A}, logrando una predicción certera de la clase '$0$'. El mapeo \textit{Twoway} y la matriz de costo acumulado (\textit{Threeway}) confirmaron la capacidad de DTW para compensar las diferencias no lineales en la velocidad de pronunciación del hablante.

\begin{figure}[H]
    \centering
    \caption{Alineación Temporal Cruzada (Twoway Mapping).}
    \includegraphics[width=0.7\textwidth]{Imagenes/results(t7).png}
    \label{fig:AdVoz3}
\end{figure}

\begin{figure}[H]
    \centering
    \caption{Ruta Óptima en Matriz de Costo Acumulado (Threeway).}
    \includegraphics[width=0.7\textwidth]{Imagenes/results(t7.1).png}
    \label{fig:AdVoz3}
\end{figure}

\section{Conclusiones}
\begin{enumerate}
    \item La limpieza adaptativa de guiones y el análisis de N-gramas permiten extraer componentes narrativos clave como la trama de la clase especial en \textit{Malcolm}.
    \item La lematización con \textit{WordNet} combinada con TF-IDF a nivel palabra ofreció el mejor rendimiento en clasificación temática de texto ($94.72\%$).
    \item La calibración del suavizado gaussiano previo a la Transformada de Hough permitió cuantificar con precisión 134 componentes industriales eliminando falsos positivos por brillo.
    \item El aprendizaje por transferencia mediante redes VGG16 congeladas demostró ser una alternativa eficiente para diagnóstico biomédico en regímenes de datos escasos.
    \item El Centroide Espectral cuantificó la distancia acústica ($1,296.46\text{ Hz}$) entre arreglos musicales dramáticos y festivos.
    \item La combinación de normalización de amplitud, filtrado Wavelet $cA_4$ y alineación temporal mediante DTW proporcionó un esquema altamente discriminativo para el reconocimiento de voz dependiente de la forma espectral ($d_{\text{$0$}} = 70.7006 < d_{\text{$1$}} = 75.4533$).
\end{enumerate}

En conjunto, la integración de estas metodologías confirma que la combinación estratégica de técnicas de procesamiento constituye una herramienta sumamente sólida y versátil. Este enfoque no solo optimiza la extracción de características altamente discriminativas a partir de fuentes de datos complejas e inestructuradas (texto, imágenes y audio), sino que también garantiza la interpretabilidad, precisión y eficiencia computacional requeridas en aplicaciones del mundo real.

\clearpage
% 7. REFERENCIAS
\nocite{*}
\bibliographystyle{unsrtnat} 
\bibliography{referencias}

\end{document}